% ID: SEL-D0008-BC-E06
% ESTADO: validada. S009, 2026-11-07; tres pasadas del mismo autor y revisión matemática, didáctica y visual.
\textbf{Elección de la urna.}\par\nopagebreak[4]
Usamos la convención habitual de dos dados de seis caras, equilibrados e independientes: hay 36 pares ordenados equiprobables. Las sumas 9,10,11,12 tienen respectivamente 4,3,2,1 resultados; por tanto,
\[
P(A)=\frac{10}{36}=\frac5{18},\qquad P(B)=\frac{26}{36}=\frac{13}{18}.
\]
No se consideran equiprobables las once sumas posibles. Sin esa convención sobre los dados harían falta sus probabilidades para obtener una respuesta numérica.

\textbf{a) Verde y de B.}\par\nopagebreak[4]
Sea $V$ ``extraer verde''. La pregunta es una intersección:
\[
P(V\cap B)=P(B)P(V\mid B)
=\frac{13}{18}\frac39=\boxed{\frac{13}{54}\approx0.240741}.
\]

\textbf{b) Probabilidad total de roja.}\par\nopagebreak[4]
Sea $R$ ``extraer roja''. Como las urnas forman una partición,
\[
P(R)=P(A)P(R\mid A)+P(B)P(R\mid B)
=\frac5{18}\frac49+\frac{13}{18}\frac69
=\boxed{\frac{49}{81}\approx0.604938}.
\]
Comprobación: $P(V)=(5/18)(5/9)+(13/18)(3/9)=32/81$, complementaria de $49/81$. Sólo se extrae una bola; no es necesario introducir reposición entre extracciones.
